% ===== PRÉAMBULE CANONIQUE — à coller VERBATIM en tête de CHAQUE .tex =====
\documentclass[11pt,a4paper]{article}
\usepackage[utf8]{inputenc}\usepackage[T1]{fontenc}\usepackage{lmodern}
\usepackage[french]{babel}
\usepackage{amsmath,amssymb,mathtools}\usepackage{icomma}
\usepackage[version=4]{mhchem}          % \ce{...} : équations & formules
\usepackage{siunitx}\sisetup{locale=FR} % \qty{}{}, \si{}, \ang{}
\usepackage{chemfig}                    % \chemfig{...} : structures développées
\usepackage{xcolor}\usepackage{colortbl}
\usepackage{tikz}\usepackage{pgfplots}\pgfplotsset{compat=1.18}
\usepackage[most]{tcolorbox}\usepackage{enumitem}\usepackage{array,booktabs}
\usepackage[a4paper,margin=2.2cm]{geometry}
\usetikzlibrary{arrows.meta,calc,angles,quotes,positioning,patterns,backgrounds,shapes.geometric}
\definecolor{primary}{RGB}{222,49,99}\definecolor{primarydark}{RGB}{150,22,55}
\definecolor{defcol}{RGB}{0,114,178}\definecolor{methcol}{RGB}{52,92,148}
\definecolor{excol}{RGB}{0,135,90}\definecolor{attncol}{RGB}{200,40,55}
\tcbset{boxbase/.style={enhanced,breakable,boxrule=0.9pt,arc=2.5pt,
  left=3mm,right=3mm,top=2.3mm,bottom=2.3mm,fonttitle=\bfseries,coltitle=white}}
% --- boîtes TUTORIEL (noms EXACTS, ne pas renommer) ---
\newtcolorbox{cle}[1][]{boxbase,colback=defcol!6,colframe=defcol,
  title={\textsc{À retenir}\ifx\relax#1\relax\relax\else\textmd{ : \itshape #1}\fi}}
\newtcolorbox{methode}[1][]{boxbase,colback=methcol!7,colframe=methcol,sharp corners,
  title={\textsc{Méthode}\ifx\relax#1\relax\relax\else\textmd{ --- \itshape #1}\fi}}
\newtcolorbox{exemplebox}[1][]{boxbase,colback=excol!5,colframe=excol,
  title={\textsc{Exemple}\ifx\relax#1\relax\relax\else\textmd{ : \itshape #1}\fi}}
\newtcolorbox{piege}[1][]{boxbase,colback=attncol!6,colframe=attncol,
  title={\textsc{Piège}\ifx\relax#1\relax\relax\else\textmd{ : \itshape #1}\fi}}
\newtcolorbox{remarque}[1][]{boxbase,colback=black!4,colframe=black!45,
  title={\textsc{Remarque}\ifx\relax#1\relax\relax\else\textmd{ : \itshape #1}\fi}}
% --- boîtes EXERCICE / CORRIGÉ (noms EXACTS, auto-comptées) ---
\newtcolorbox[auto counter]{exo}[1][]{boxbase,colback=primary!4,colframe=primary,
  title={\textsc{Exercice}~\thetcbcounter\ifx\relax#1\relax\relax\else\textmd{ --- \itshape #1}\fi}}
\newtcolorbox[auto counter]{corrige}[1][]{boxbase,colback=excol!4,colframe=excol,
  title={\textsc{Corrigé}~\thetcbcounter\ifx\relax#1\relax\relax\else\textmd{ --- \itshape #1}\fi}}
\newcommand{\R}{\mathbb{R}}\newcommand{\N}{\mathbb{N}}\newcommand{\Z}{\mathbb{Z}}\newcommand{\ds}{\displaystyle}
% (un \usepackage{fancyhdr}+en-tête est possible ; il est ignoré côté web)
% =========================================================================
\begin{document}

\begin{center}{\large\bfseries\color{primarydark} Thème 2 — Niveau Facile}\end{center}
\emph{Données :} $\rho_{\text{eau}} = \SI{1.00}{\gram\per\milli\litre}$ ; aux CNTP
$V_m = \SI{22.4}{\litre\per\mole}$.

\begin{exo}[calculer une masse volumique]
Détermine la masse volumique $\rho$ ($\rho = m/V$).
\begin{enumerate}[label=\alph*)]
  \item $\SI{24}{\gram}$ occupant $\SI{30}{\milli\litre}$ ;
  \item $\SI{79}{\gram}$ occupant $\SI{100}{\milli\litre}$ ;
  \item $\SI{13.55}{\gram}$ occupant $\SI{1}{\milli\litre}$ (mercure).
\end{enumerate}
\end{exo}

\begin{exo}[de la masse volumique à la masse]
Quelle masse représente le volume indiqué ($m = \rho\,V$)~? \emph{Attention aux unités
($\SI{1}{\litre} = \SI{1000}{\milli\litre}$).}
\begin{enumerate}[label=\alph*)]
  \item $\SI{1}{\litre}$ d'une solution de $\rho = \SI{1.10}{\gram\per\milli\litre}$ ;
  \item $\SI{250}{\milli\litre}$ d'un liquide de $\rho = \SI{0.80}{\gram\per\milli\litre}$ ;
  \item $\SI{2}{\litre}$ d'une solution de $\rho = \SI{1.20}{\gram\per\milli\litre}$.
\end{enumerate}
\end{exo}

\begin{exo}[densité et masse volumique]
\begin{enumerate}[label=\alph*)]
  \item L'éthanol a une densité $d = 0{,}79$. Quelle est sa masse volumique~?
  \item Le mercure a $\rho = \SI{13.6}{\gram\per\milli\litre}$. Quelle est sa densité~?
  \item Une solution a $d = 1{,}03$. Quelle est sa masse volumique en \si{\gram\per\milli\litre}~?
\end{enumerate}
\end{exo}

\begin{exo}[volume d'un gaz aux CNTP]
Quel volume occupe, aux CNTP, la quantité de gaz indiquée ($V = n\,V_m$)~?
\begin{enumerate}[label=\alph*)]
  \item $\SI{1}{\mole}$ ; \quad b) $\SI{0.5}{\mole}$ ; \quad c) $\SI{0.25}{\mole}$ ; \quad d) $\SI{2}{\mole}$.
\end{enumerate}
\end{exo}

\begin{exo}[du volume de gaz à la quantité de matière]
Quelle quantité de matière représente, aux CNTP ($n = V/V_m$)~?
\begin{enumerate}[label=\alph*)]
  \item $\SI{22.4}{\litre}$ ; \quad b) $\SI{11.2}{\litre}$ ; \quad c) $\SI{44.8}{\litre}$ ; \quad d) $\SI{5.6}{\litre}$.
\end{enumerate}
\end{exo}

\end{document}
